% references.tex
\chapter{Comprehensive Bibliography}
\addcontentsline{toc}{chapter}{Bibliography}

\begin{itemize}[leftmargin=1.5cm]
    \item Albergo, M. S., Kanwar, G., \& Shanahan, P. E. (2019). Flow-based generative models for Markov chain Monte Carlo in lattice field theory. \textit{Physical Review D}, 100(3), 034515.
    \item Bellemare, M. G., et al. (2020). Autonomous navigation of stratospheric balloons using reinforcement learning. \textit{Nature}, 588(7836), 77-82.
    \item Carleo, G., \& Troyer, M. (2017). Solving the quantum many-body problem with artificial neural networks. \textit{Science}, 355(6325), 602-606.
    \item Carleo, G., et al. (2019). Machine learning and the physical sciences. \textit{Reviews of Modern Physics}, 91(4), 045002.
    \item Carrasquilla, J., \& Melko, R. G. (2017). Machine learning phases of matter. \textit{Nature Physics}, 13(5), 431-434.
    \item Cranmer, M., et al. (2020). Lagrangian Neural Networks. \textit{arXiv preprint arXiv:2003.04638}.
    \item Degrave, J., et al. (2022). Magnetic control of tokamak plasmas through deep reinforcement learning. \textit{Nature}, 602(7897), 414-419.
    \item Duarte, J., et al. (2018). Fast inference of deep neural networks in FPGAs for particle physics. \textit{Journal of Instrumentation}, 13(07), P07027.
    \item Gabbard, H., et al. (2018). Matching filtered search for gravitational waves from binary black holes using deep learning. \textit{Physical Review Letters}, 120(14), 141103.
    \item George, D., \& Huerta, E. A. (2018). Deep learning for real-time gravitational wave detection and parameter estimation: Results with advanced LIGO data. \textit{Physics Letters B}, 778, 64-70.
    \item Gomez-Bombarelli, R., et al. (2018). Automatic chemical design using a data-driven continuous representation of molecules. \textit{ACS Central Science}, 4(2), 268-276.
    \item Goodfellow, I., Bengio, Y., \& Courville, A. (2016). \textit{Deep Learning}. MIT Press.
    \item Greydanus, S., Dzamba, M., \& Yosinski, J. (2019). Hamiltonian Neural Networks. \textit{Advances in Neural Information Processing Systems}, 32.
    \item Harris, C. R., et al. (2020). Array programming with NumPy. \textit{Nature}, 585(7825), 357-362.
    \item Karniadakis, G. E., et al. (2021). Physics-informed machine learning. \textit{Nature Reviews Physics}, 3(6), 422-440.
    \item Mehta, P., et al. (2019). A high-bias, low-variance introduction to Machine Learning for physicists. \textit{Physics Reports}, 810, 1-124.
    \item Noé, F., et al. (2019). Boltzmann generators: Sampling equilibrium states of many-body systems with deep learning. \textit{Science}, 365(6457).
    \item Raissi, M., Perdikaris, P., \& Karniadakis, G. E. (2019). Physics-informed neural networks: A deep learning framework for solving forward and inverse problems involving nonlinear partial differential equations. \textit{Journal of Computational Physics}, 378, 686-707.
    \item Udrescu, S. M., \& Tegmark, M. (2020). AI Feynman: A physics-inspired symbolic regressor. \textit{Science Advances}, 6(16), eaay2631.
    \item Virtanen, P., et al. (2020). SciPy 1.0: fundamental algorithms for scientific computing in Python. \textit{Nature Methods}, 17(3), 261-272.
\end{itemize}
